% !TEX root = ../thesis.tex
% ==============================================================================
%  APPENDIX: 
%       DICTIONARY & GLOSSARY PIPELINE
%       (appendices/appendix_dictionary_and_glossary.tex)
% ==============================================================================
%  This appendix documents the automated bidirectional glossary system,
%  Unicode Persian collation rules, and compiler integration pipelines.
% ==============================================================================


\section{سازوکار واژه‌نامه دوزبانه و مرتب‌سازی الفبایی خودکار}\label{sec:app_dictionary_and_glossary}

ارائه‌ی واژه‌نامه‌های تخصصی در انتهای پایان‌نامه به دو صورت «فارسی به انگلیسی» و «انگلیسی به فارسی»، از الزامات رایج دانشگاه‌هاست. در روش‌های سنتی، پژوهشگر ناچار بود واژه‌ها را در دو فایل جداگانه به‌صورت دستی الفبایی نماید که فرآیندی زمان‌بر و مستعد خطا بود.

در این قالب، سامانه‌ی واژه‌نامه بر اساس مفهوم «مبدأ واحد حقیقت» پیاده‌سازی شده است؛ بدین صورت که هر دو واژه‌نامه به صورت اتوماتیک بر اساس یک فایل ساخته می‌شوند. لذا تمام واژه‌های تخصصی در یک فایل مرکزی به نشانی \code{backmatter/words.tex} با کمک دستور زیر ثبت می‌شوند:
\begin{center}
\code{\textbackslash{}dictentry\{English Term\}\{\rl{معادل فارسی}\}}
\end{center}


\subsection{خط‌لوله‌ی مرتب‌سازی خودکار بر مبنای یونیکد}

اسکریپت پایتون \code{scripts/build\_dictionaries.py} وظیفه‌ی خواندن فایل \code{words.tex}، مرتب‌سازی واژه‌ها و تولید دو فایل \code{dict\_fa2en\_content.tex} و \code{dict\_en2fa\_content.tex} را بر عهده دارد. اصول حاکم بر این اسکریپت عبارتند از:
\begin{itemize}
    \item \textbf{رعایت دقیق ترتیب الفبای فارسی}: مرتب‌سازی بر اساس حروف ۳۲‌-گانه‌ی فارسی از «آ» و «الف» تا «ی» با لحاظ کردن جایگاه صحیح حروف خاص (گ، چ، پ، ژ) انجام می‌پذیرد.
    \item \textbf{پالایش نیم‌فاصله‌ها در کلید مرتب‌سازی}: برای جلوگیری از اختلال در چیدمان الفبایی، کاراکترهای نیم‌فاصله (\lr{ZWNJ} با کد یونیکد \code{{\textbackslash}u200c}) و خط‌های تیره در زمان تعیین رتبه‌ی الفبایی نادیده گرفته می‌شوند؛ بنابراین واژه‌هایی مانند «شبکه‌ی عصبی» و «شبکه عصبی» دقیقاً در جایگاه الفبایی صحیح خود قرار می‌گیرند.
    \item \textbf{سربرگ‌های الفبایی تفکیک‌شده}: واژه‌ها به‌طور خودکار بر اساس حرف نخست گروه‌بندی شده و ماکروی \code{{\textbackslash}dicalphabet\{\ldots\}} پیش از هر دسته درج می‌گردد. همچنین برای جلوگیری از قرارگیری سربرگ تنها در انتهای ستون، تمهیدات لازم در ماکروهای مربوطه اندیشیده شده است.
\end{itemize}


\subsection{یکپارچه‌سازی با فرآیند کامپایل (\lr{latexmkrc})}

برای بی‌نیازی از اجرای دستی اسکریپت پس از افزودن هر واژه، در فایل پیکربندی \code{latexmkrc} یک قلاب کامپایل تعریف شده است. بدین ترتیب، با هر بار اجرای کامپایلر برای تولید \lr{PDF} (چه از طریق افزونه‌ی \lr{LaTeX Workshop} در \lr{VS Code} و چه از طریق خط فرمان \lr{latexmk})، اسکریپت پایتون پیش از شروع فرآیند حروف‌چینی به‌صورت خودکار اجرا شده و واژه‌نامه‌ها را به‌روزرسانی می‌کند.
